% !TeX root = ../main.tex
\section{Воздействие сигнала и шума на ЧД}

\subsection{Решение задачи 3-1}

На входе приёмника ЧМ сигналов действует немодулированное колебание и белый шум, причём $a \gg 3$. Как изменится среднеквадратическое (эффективное) напряжение шума на выходе БНЧ приёмника, если:
\begin{enumerate}
    \item полосу пропускания БВЧ увеличить вдвое, сохранив коэффициент усиления? 
    \item вдвое расширить полосу пропускания БНЧ? 
\end{enumerate}
АЧХ БВЧ и БНЧ имеют прямоугольную форму.

\medskip
\noindent\textbf{Решение:}

Так как по условию $a \gg 3$, приёмник работает в надпороговой области, и для расчёта спектральной плотности мощности шума на выходе ЧД используется выражение:
\begin{equation}
    G_{\text{ЧД}}(f) \approx S_{\text{ЧД}}^2 \frac{f^2}{U_c^2} \cdot 2 K_{0\,\text{БВЧ}}^2 G_0.
\end{equation}

Мощность (дисперсия) шума на выходе БНЧ представляет собой интеграл от спектральной плотности в пределах полосы пропускания БНЧ $[0; F_{\text{в}}]$:
\begin{equation}
\begin{aligned}
    U_{\text{ш.вых}}^2 &= K_{\text{БНЧ}}^2 \int_0^{F_{\text{в}}} G_{\text{ЧД}}(f)\,df = K_{\text{БНЧ}}^2 \int_0^{F_{\text{в}}} S_{\text{ЧД}}^2 \frac{f^2}{U_c^2} \cdot 2 K_{0\,\text{БВЧ}}^2 G_0\,df = \\
    &= \frac{2 K_{\text{БНЧ}}^2 S_{\text{ЧД}}^2 K_{0\,\text{БВЧ}}^2 G_0}{U_c^2} \cdot \frac{F_{\text{в}}^3}{3}.
\end{aligned}
\end{equation}

    Согласно приведённой формуле спектральной плотности мощности шума, $U_{\text{ш.вых}}^2$ зависит от полосы БНЧ и не зависит от полосы пропускания БВЧ. При изменении полосы пропускания БВЧ в 2 раза с сохранением коэффициента усиления эффективное напряжение шума на выходе БНЧ не изменится.
    
    Если вдвое расширить полосу пропускания БНЧ, верхняя граничная частота увеличится в 2 раза ($F_{\text{в}}' = 2 F_{\text{в}}$). Мощность шума на выходе зависит от куба этой частоты, следовательно, $U_{\text{ш.вых}}^2$ увеличится в $2^3 = 8$ раз. Эффективное значение напряжения шума на выходе БНЧ увеличится в $\sqrt{8} \approx 2{,}83$ раза.

\noindent\textbf{Ответ:} а) не изменится; б) увеличится в $2{,}83$ раза.

\subsection{Решение задачи 3-2}
